\documentclass[12pt]{article}
\usepackage[utf8]{inputenc}
\usepackage[french]{babel}
\usepackage{amsmath}
\usepackage[many]{tcolorbox}
\usepackage[siunitx]{circuitikz}
\usepackage{pgfplots}
\usepackage{chemfig}
\usetikzlibrary{arrows.meta}
\usetikzlibrary{babel}
\usetikzlibrary{shapes.symbols}
\usetikzlibrary{patterns}
\usetikzlibrary{calc}

\begin{document}

\section*{Circuit RC et Charge d'un Condensateur}

\subsection*{Schéma du Circuit RC}
Voici le schéma d'un circuit RC en série :
\begin{center}
\begin{tikzpicture}[american, font=\sffamily]
    \draw
    (0,0) node[resistor, draw, name=r] {$R$}
    -- ++(0,3) node[ground, rotate=-90, draw, name=gnd] {\;}
    -- ++(0,0) node[battery1, draw, name=vs, rotate=-90, voltage={$\mathcal{V}_0$}] {}
    -- ++(-0.5,0) node[ocirc]{}
    -- ++(0,-3) node[ocirc]{}
    -- ++(0,-0.5) node[cspst, draw, name=s] {}
    -- ++(0,-1) node[ocirc]{}
    -- ++(2,0)
    -- ++(0,4) node[capacitor, draw, name=c] {$C$}
    -- ++(-2,0)
    -- (c.west) node[opamp symbol, draw, rotate=90, yscale=-1, name=opamp] {}
    -- (opamp.east);
    \draw (vs.east) -- ++(0.5,0) node[above] {$\mathcal{V}_0$} -- (vs.east |- c.north);
\end{tikzpicture}
\end{center}

\subsection*{Graphique de la Charge du Condensateur}
Voici la courbe de charge d'un condensateur dans un circuit RC en fonction du temps :

\begin{center}
\begin{tikzpicture}
\begin{axis}[
    xlabel={Temps $t$ ($\si{\second}$)},
    ylabel={Tension aux bornes du condensateur $V_C(t)$ (\si{\volt})},
    xmin=0, xmax=0.25,
    ymin=0, ymax=10,
    samples=300,
    axis lines=middle,
    grid=major,
    width=10cm,
    height=6cm,
    style={thick}
]

% Paramètres
\def\Vzero{10}
\def\RCtime{0.1}

\addplot[blue, domain=0:0.25, samples=300] {\Vzero * (1 - exp(-x/\RCtime))};

\addlegendentry{$\mathcal{V}_C(t) = \mathcal{V}_0 (1 - \exp(-t/RC))$}

\draw[dashed, red] (\RCtime,0) -- (\RCtime, {\Vzero * (1 - exp(-1))});
\draw[dashed, red] (0,{\Vzero * (1 - exp(-1))}) -- (\RCtime,{\Vzero * (1 - exp(-1))});

\node[above] at (axis cs:\RCtime,{\Vzero * (1 - exp(-1))}) {$\mathcal{V}_C(t) = \mathcal{V}_0 (1 - \frac{1}{e})$};
\node[below] at (axis cs:\RCtime,0) {$t_{RC}$};
\end{axis}
\end{tikzpicture}

*Équation théorique* :
\[
V_C(t) = \mathcal{V}_0 \left(1 - \exp\left(-\frac{t}{RC}\right)\right)
\]
\end{center}

\end{document}